\chapter{1987年广东卷（文）}

\section{单选题}

\begin{problem}
\(\cos1030^{\circ}=\)（\quad）

\choices
{\(\cos50^{\circ}\)}
{\(-\cos50^{\circ}\)}
{\(\sin50^{\circ}\)}
{\(-\sin50^{\circ}\)}
\end{problem}

\begin{answer}A\end{answer}

\begin{solution}
因为 \(1030^{\circ}=3\cdot360^{\circ}-50^{\circ}\)，所以 \(\cos1030^{\circ}=\cos(-50^{\circ})=\cos50^{\circ}\)，故选 A.
\end{solution}

\begin{problem}
\(2^{6}\) 与 \(0.5^{2}\) 的等比中项是（\quad）

\choices
{\(16\)}
{\(\pm4\)}
{\(2\)}
{\(4\)}
\end{problem}

\begin{answer}B\end{answer}

\begin{solution}
设等比中项为 \(m\)，则 \(m^{2}=2^{6}\cdot0.5^{2}=64\cdot\frac14=16\)，所以 \(m=\pm4\)，故选 B.
\end{solution}

\begin{problem}
函数 \(y=\sqrt{2x-x^{2}}\) 的定义域是（\quad）

\choices
{\((-\infty,0)\)}
{\((0,2]\)}
{\([0,2]\)}
{\([-2,0]\)}
\end{problem}

\begin{answer}C\end{answer}

\begin{solution}
根式有意义的条件是 \(2x-x^{2}\geqslant0\)，即 \(x(2-x)\geqslant0\). 因此 \(0\leqslant x\leqslant2\)，定义域为 \([0,2]\)，故选 C.
\end{solution}

\begin{problem}
指数方程 \(2\cdot4^{x}-5\cdot2^{x}+2=0\) 的解集是（\quad）

\choices
{\(\{1\}\)}
{\(\left\{\frac12,1\right\}\)}
{\(\left\{-1,\frac12\right\}\)}
{\(\{-1,1\}\)}
\end{problem}

\begin{answer}D\end{answer}

\begin{solution}
令 \(t=2^{x}>0\)，则 \(4^{x}=t^{2}\)，原方程化为 \(2t^{2}-5t+2=0\)，即 \((2t-1)(t-2)=0\). 所以 \(t=\frac12\) 或 \(t=2\)，从而 \(x=-1\) 或 \(x=1\)，解集为 \(\{-1,1\}\)，故选 D.
\end{solution}

\begin{problem}
下列不等式中，成立的是（\quad）

\choices
{\(\sin\frac{9\pi}{7}>\sin\frac{8\pi}{7}\)}
{\(\cos\frac{9\pi}{7}>\cos\frac{8\pi}{7}\)}
{\(\sec\frac{9\pi}{7}>\sec\frac{8\pi}{7}\)}
{\(\cot\frac{9\pi}{7}>\cot\frac{8\pi}{7}\)}
\end{problem}

\begin{answer}B\end{answer}

\begin{solution}
因为 \(\frac{8\pi}{7}=\pi+\frac\pi7\)，\(\frac{9\pi}{7}=\pi+\frac{2\pi}{7}\)，所以 \(\sin\frac{8\pi}{7}=-\sin\frac\pi7\)，\(\sin\frac{9\pi}{7}=-\sin\frac{2\pi}{7}\). 在 \(\left(0,\frac\pi2\right)\) 上正弦递增，故选项 A 不成立.

由于 \(0<\frac\pi7<\frac{2\pi}{7}<\frac\pi2\)，且余弦在 \(\left(0,\frac\pi2\right)\) 上递减，有 \(\cos\frac{2\pi}{7}<\cos\frac\pi7\). 因而
\(\cos\frac{9\pi}{7}=-\cos\frac{2\pi}{7}>-\cos\frac\pi7=\cos\frac{8\pi}{7}\)，选项 B 成立. 两个余弦均为负数，取倒数后不等号反向，所以 \(\sec\frac{9\pi}{7}<\sec\frac{8\pi}{7}\)，选项 C 不成立. 又 \(\frac{8\pi}{7}\)，\(\frac{9\pi}{7}\in\left(\pi,\frac{3\pi}{2}\right)\)，余切在该区间上递减，故 \(\cot\frac{9\pi}{7}<\cot\frac{8\pi}{7}\)，选项 D 不成立. 因此选 B.
\end{solution}

\begin{problem}
如果圆锥的高和底面直径都等于 \(a\)，则该圆锥的体积是（\quad）

\choices
{\(\frac{\pi}{4}a^{3}\)}
{\(\frac{\pi}{6}a^{3}\)}
{\(\frac{\pi}{3}a^{3}\)}
{\(\frac{\pi}{12}a^{3}\)}
\end{problem}

\begin{answer}D\end{answer}

\begin{solution}
圆锥底面直径为 \(a\)，所以半径为 \(\frac a2\)，高为 \(a\). 因此体积 \(V=\frac13\pi\left(\frac a2\right)^{2}a=\frac\pi{12}a^{3}\)，故选 D.
\end{solution}

\begin{problem}
在圆 \(x^{2}+y^{2}=4\) 上，与直线 \(4x+3y-12=0\) 的距离最小的点的坐标是（\quad）

\choices
{\(\left(\frac85,\frac65\right)\)}
{\(\left(\frac85,-\frac65\right)\)}
{\(\left(-\frac85,\frac65\right)\)}
{\(\left(-\frac85,-\frac65\right)\)}
\end{problem}

\begin{answer}A\end{answer}

\begin{solution}
圆 \(x^{2}+y^{2}=4\) 的半径为 \(2\). 对圆上任意点，有 \(4x+3y\leqslant\sqrt{4^{2}+3^{2}}\sqrt{x^{2}+y^{2}}=10<12\)，故点到直线的距离为 \(\frac{12-4x-3y}{5}\). 要使距离最小，只需使 \(4x+3y\) 最大.

由柯西不等式，等号在 \((x,y)\) 与 \((4,3)\) 同方向时取得，结合 \(x^{2}+y^{2}=4\)，得 \((x,y)=2\left(\frac45,\frac35\right)=\left(\frac85,\frac65\right)\). 故选 A.
\end{solution}

\begin{problem}
若复数 \(z\) 的辐角是 \(\frac56\pi\)，实部是 \(-2\sqrt3\)，则 \(z=\)（\quad）

\choices
{\(-2\sqrt3-2\i\)}
{\(-2\sqrt3+2\i\)}
{\(-2\sqrt3+2\sqrt3\i\)}
{\(-2\sqrt3-2\sqrt3\i\)}
\end{problem}

\begin{answer}B\end{answer}

\begin{solution}
设 \(z=r\left(\cos\frac{5\pi}{6}+\i\sin\frac{5\pi}{6}\right)\). 由实部条件，\(r\cos\frac{5\pi}{6}=r\left(-\frac{\sqrt3}{2}\right)=-2\sqrt3\)，所以 \(r=4\). 于是虚部为 \(4\sin\frac{5\pi}{6}=2\)，从而 \(z=-2\sqrt3+2\i\)，故选 B.
\end{solution}

\begin{problem}
已知函数 \(f(x)=x^{\frac13}\)，\(g(x)=\left(\frac13\right)^x\)，那么在 \((-\infty,+\infty)\) 上（\quad）

\choices
{\(f(x)\) 和 \(g(x)\) 都是增函数}
{\(f(x)\) 和 \(g(x)\) 都是减函数}
{\(f(x)\) 是减函数，而 \(g(x)\) 是增函数}
{\(f(x)\) 是增函数，而 \(g(x)\) 是减函数}
\end{problem}

\begin{answer}D\end{answer}

\begin{solution}
若 \(x_{1}<x_{2}\)，由于立方函数严格递增，有 \(x_{1}^{1/3}<x_{2}^{1/3}\)，所以 \(f(x)=x^{1/3}\) 在 \(\R\) 上是增函数. 又因为 \(0<\frac13<1\)，指数函数 \(g(x)=\left(\frac13\right)^{x}\) 在 \(\R\) 上是减函数. 因此选 D.
\end{solution}

\begin{problem}
\(\displaystyle\frac{(\sin\alpha+\cos\alpha)^{2}-(\sin\alpha-\cos\alpha)^{2}}{\tan\alpha-\sin\alpha\cos\alpha}=\)（\quad）

\choices
{\(2\cot^{2}\alpha\)}
{\(2\tan^{2}\alpha\)}
{\(4\cot^{2}\alpha\)}
{\(4\tan^{2}\alpha\)}
\end{problem}

\begin{answer}C\end{answer}

\begin{solution}
分子为 \((\sin\alpha+\cos\alpha)^{2}-(\sin\alpha-\cos\alpha)^{2}=4\sin\alpha\cos\alpha\). 在原式有意义时，\(\tan\alpha-\sin\alpha\cos\alpha=\frac{\sin\alpha}{\cos\alpha}-\sin\alpha\cos\alpha=\frac{\sin^{3}\alpha}{\cos\alpha}\). 所以原式 \(=\frac{4\sin\alpha\cos\alpha}{\sin^{3}\alpha/\cos\alpha}=4\cot^{2}\alpha\)，故选 C.
\end{solution}

\begin{problem}
函数 \(y=\sin3x+2\cos3x\) 的最小值是（\quad）

\choices
{\(-2\)}
{\(-\sqrt5\)}
{\(-3\)}
{\(-5\)}
\end{problem}

\begin{answer}B\end{answer}

\begin{solution}
由柯西不等式，\(\sin3x+2\cos3x\geqslant-\sqrt{1^{2}+2^{2}}=-\sqrt5\). 取 \(\phi\) 满足 \(\cos\phi=\frac1{\sqrt5}\)、\(\sin\phi=\frac2{\sqrt5}\)，则
\(\sin3x+2\cos3x=\sqrt5\sin(3x+\phi)\). 取 \(3x+\phi=-\frac\pi2+2k\pi\) 即可取得 \(-\sqrt5\)，所以最小值为 \(-\sqrt5\)，故选 B.
\end{solution}

\begin{problem}
若 \(\cot\frac{\theta}{2}=3\)，则 \(\cos\theta\) 的值是（\quad）

\choices
{\(\frac35\)}
{\(-\frac35\)}
{\(-\frac45\)}
{\(\frac45\)}
\end{problem}

\begin{answer}D\end{answer}

\begin{solution}
由 \(\cot\frac\theta2=3\)，得 \(\tan\frac\theta2=\frac13\). 利用半角正切公式，\(\cos\theta=\frac{1-\tan^{2}\frac\theta2}{1+\tan^{2}\frac\theta2}=\frac{1-\frac19}{1+\frac19}=\frac45\)，故选 D.
\end{solution}

\begin{problem}
不等式 \(\displaystyle\frac{3x-1}{2-x}\geqslant1\) 的解集是（\quad）

\choices
{\(\left\{x\mid\frac34\leqslant x\leqslant2\right\}\)}
{\(\left\{x\mid\frac34\leqslant x<2\right\}\)}
{\(\left\{x\mid x>2\text{或}x\leqslant\frac34\right\}\)}
{\(\{x\mid x<2\}\)}
\end{problem}

\begin{answer}B\end{answer}

\begin{solution}
移项得 \(\frac{3x-1}{2-x}-1=\frac{4x-3}{2-x}\geqslant0\)，且 \(x\neq2\). 临界点为 \(\frac34\) 和 \(2\). 当 \(x<\frac34\) 时，\(4x-3<0\)、\(2-x>0\)，分式为负；当 \(\frac34\leqslant x<2\) 时，\(4x-3\geqslant0\)、\(2-x>0\)，分式为非负；当 \(x>2\) 时，\(4x-3>0\)、\(2-x<0\)，分式为负. 所以解集为 \(\left[\frac34,2\right)\)，选 B.
\end{solution}

\begin{problem}
设 \(f(x)=\frac{2x+1}{4x+3}\)（\(x\in\R\)，\(x\neq-\frac34\)），则 \(f^{-1}(2)=\)（\quad）

\choices
{\(-\frac56\)}
{\(\frac5{11}\)}
{\(\frac25\)}
{\(-\frac25\)}
\end{problem}

\begin{answer}A\end{answer}

\begin{solution}
令 \(f(x)=2\)，则 \(\frac{2x+1}{4x+3}=2\)，从而 \(2x+1=8x+6\)，解得 \(x=-\frac56\). 此时 \(4x+3=-\frac13\neq0\)，所以 \(f^{-1}(2)=-\frac56\)，故选 A.
\end{solution}

\begin{problem}
用 \(0\)，\(1\)，\(2\)，\(3\) 这四个数字，组成个位数不是 \(1\) 的没有重复数字的四位数，共有（\quad）

\choices
{\(16\) 个}
{\(14\) 个}
{\(12\) 个}
{\(10\) 个}
\end{problem}

\begin{answer}B\end{answer}

\begin{solution}
四位数的个位不能为 \(1\). 若个位为 \(0\)，千位可从 \(1\)，\(2\)，\(3\) 中选，有 \(3\) 种，其余两位有 \(2\) 种排列，共 \(6\) 个；若个位为 \(2\)，千位不能为 \(0\)，只能从 \(1,3\) 中选，有 \(2\) 种，其余两位有 \(2\) 种排列，共 \(4\) 个；若个位为 \(3\)，千位同样只能从 \(1,2\) 中选，有 \(2\) 种，其余两位有 \(2\) 种排列，共 \(4\) 个. 因此共有 \(6+4+4=14\) 个，故选 B.
\end{solution}

\begin{problem}
抛物线 \(y=x^{2}+3x-2\) 的对称轴方程是（\quad）

\choices
{\(x=\sqrt3\)}
{\(x=-\sqrt3\)}
{\(x=\frac32\)}
{\(x=-\frac32\)}
\end{problem}

\begin{answer}D\end{answer}

\begin{solution}
将 \(y=x^{2}+3x-2\) 写成 \(y=\left(x+\frac32\right)^{2}-\frac{17}{4}\). 抛物线的对称轴为 \(x=-\frac32\)，故选 D.
\end{solution}

\begin{problem}
设 \(A\)，\(B\) 是直线 \(3x+4y+2=0\) 与圆 \(x^{2}+y^{2}+4y=0\) 的两个交点，则线段 \(AB\) 的垂直平分线方程是（\quad）

\choices
{\(3y-4x+2=0\)}
{\(3y-4x+6=0\)}
{\(3y+4x+6=0\)}
{\(4y+3x+8=0\)}
\end{problem}

\begin{answer}B\end{answer}

\begin{solution}
圆 \(x^{2}+y^{2}+4y=0\) 的圆心为 \(O(0,-2)\). 弦 \(AB\) 在直线 \(3x+4y+2=0\) 上，所以线段 \(AB\) 的垂直平分线经过 \(O\)，且垂直于该直线. 原直线斜率为 \(-\frac34\)，垂直平分线斜率为 \(\frac43\)，故其方程为 \(y+2=\frac43x\)，即 \(3y-4x+6=0\)，故选 B.
\end{solution}

\begin{problem}
椭圆 \(x^{2}+4y^{2}+4x=0\) 的右焦点坐标是（\quad）

\choices
{\((\sqrt3-2,0)\)}
{\((\sqrt3+2,0)\)}
{\((\sqrt3,0)\)}
{\((-\sqrt3,0)\)}
\end{problem}

\begin{answer}A\end{answer}

\begin{solution}
配方得 \(x^{2}+4y^{2}+4x=0\iff\frac{(x+2)^{2}}4+y^{2}=1\). 因此椭圆中心为 \((-2,0)\)，半长轴 \(a=2\)，半短轴 \(b=1\)，焦半距 \(c=\sqrt{a^{2}-b^{2}}=\sqrt3\). 右焦点为 \((-2+\sqrt3,0)=(\sqrt3-2,0)\)，故选 A.
\end{solution}

\begin{problem}
设双曲线的半焦距为 \(c\)，两条准线的距离为 \(d\)，且 \(c=d\)，则这双曲线的离心率 \(e=\)（\quad）

\choices
{\(\sqrt3\)}
{\(\sqrt2\)}
{\(2\)}
{\(3\)}
\end{problem}

\begin{answer}B\end{answer}

\begin{solution}
设双曲线的半实轴为 \(a\)，则半焦距 \(c=ae\)，两条准线为 \(x=\pm\frac ae\)（或 \(y=\pm\frac ae\)），所以两准线距离 \(d=\frac{2a}{e}\). 由 \(c=d\) 得 \(ae=\frac{2a}{e}\)，从而 \(e^{2}=2\). 因为双曲线的离心率 \(e>1\)，所以 \(e=\sqrt2\)，故选 B.
\end{solution}

\begin{problem}
若 \(\theta\in\left(\frac\pi2,\frac34\pi\right)\)，则 \(\displaystyle\frac{x^{2}}{\csc\theta-2}+\frac{y^{2}}{2-\sec\theta}=1\) 所表示的曲线是（\quad）

\choices
{焦点在 \(x\) 轴上的椭圆}
{焦点在 \(y\) 轴上的椭圆}
{焦点在 \(x\) 轴上的双曲线}
{焦点在 \(y\) 轴上的双曲线}
\end{problem}

\begin{answer}D\end{answer}

\begin{solution}
由 \(\theta\in\left(\frac\pi2,\frac{3\pi}{4}\right)\)，得 \(1<\csc\theta<\sqrt2\)，且 \(\sec\theta<-\sqrt2\). 因此 \(\csc\theta-2<0\)、\(2-\sec\theta>0\)，并且 \(2-\csc\theta>0\). 原方程等价于
\(\displaystyle\frac{y^{2}}{2-\sec\theta}-\frac{x^{2}}{2-\csc\theta}=1\).
这是实轴在 \(y\) 轴上的双曲线，焦点在 \(y\) 轴上，故选 D.
\end{solution}

\begin{problem}
若正三棱锥的一个侧面的面积与底面的面积的比等于 \(\frac23\)，则这个三棱锥的侧面和底面所成的二面角是（\quad）

\choices
{\(15^{\circ}\)}
{\(30^{\circ}\)}
{\(45^{\circ}\)}
{\(60^{\circ}\)}
\end{problem}

\begin{answer}D\end{answer}

\begin{solution}
设正三棱锥底面边长为 \(a\)，侧面三角形的高为 \(l\). 底面面积为 \(\frac{\sqrt3}{4}a^{2}\)，一个侧面面积为 \(\frac12al\). 由面积比得 \(\frac{\frac12al}{\frac{\sqrt3}{4}a^{2}}=\frac23\)，所以 \(l=\frac{a\sqrt3}{3}\).

取底面一条棱的中点 \(D\)、底面中心 \(O\) 和棱锥顶点 \(P\). 则 \(OD=\frac{a\sqrt3}{6}\)，\(PD=l=\frac{a\sqrt3}{3}\)，且 \(PO\perp\) 底面. 侧面与底面所成的二面角的平面角为 \(\angle PDO\)，故 \(\cos\angle PDO=\frac{OD}{PD}=\frac12\)，从而二面角为 \(60^{\circ}\)，故选 D.
\end{solution}

\begin{problem}
设 \(\alpha\) 和 \(\beta\) 是两个不重合的平面，\(l\) 和 \(m\) 是两条不重合的直线，则 \(\alpha\parallel\beta\) 的一个充分条件是（\quad）

\choices
{\(l\subset\alpha\)，\(m\subset\alpha\)，且 \(l\parallel\beta\)，\(m\parallel\beta\)}
{\(l\subset\alpha\)，\(m\subset\beta\)，且 \(l\parallel m\)}
{\(l\perp\alpha\)，\(m\perp\beta\)，且 \(l\parallel m\)}
{\(l\parallel\alpha\)，\(m\parallel\beta\)，且 \(l\parallel m\)}
\end{problem}

\begin{answer}C\end{answer}

\begin{solution}
若 \(l\perp\alpha\)，\(m\perp\beta\)，且 \(l\parallel m\)，则平面 \(\alpha\) 与 \(\beta\) 的法线方向相同. 由于两平面不重合，它们必平行. 因此选项 C 是 \(\alpha\parallel\beta\) 的充分条件，故选 C.
\end{solution}

\begin{problem}
如果实数 \(a\)，\(b\)，\(c\) 不全为零，那么（\quad）

\choices
{实数 \(a\)，\(b\)，\(c\) 中，至少有一个正数}
{实数 \(a\)，\(b\)，\(c\) 中，既有正数又有负数}
{实数 \(a\)，\(b\)，\(c\) 的乘积不是零}
{实数 \(a\)，\(b\)，\(c\) 中，至多有两个零}
\end{problem}

\begin{answer}D\end{answer}

\begin{solution}
若 \(a\)，\(b\)，\(c\) 中有三个零数，则 \(a=b=c=0\)，与“不全为零”矛盾. 因此其中的零数至多有两个，故选 D.
\end{solution}

\begin{problem}
已知函数 \(f(x)=\sin\left(\frac34x+\frac52\pi\right)\)（\(x\in\R\)），\(g(x)=\lg\left(\sqrt{1+x^{2}}-x\right)\)（\(x\in\R\)），则（\quad）

\choices
{\(f(x)\) 是奇函数，而 \(g(x)\) 既不是奇函数，又不是偶函数}
{\(f(x)\) 是偶函数，而 \(g(x)\) 是奇函数}
{\(f(x)\) 和 \(g(x)\) 都是奇函数}
{\(f(x)\) 是偶函数，而 \(g(x)\) 既不是奇函数，又不是偶函数}
\end{problem}

\begin{answer}B\end{answer}

\begin{solution}
因为 \(\frac52\pi=2\pi+\frac\pi2\)，所以 \(f(x)=\cos\frac{3x}{4}\)，从而 \(f(-x)=f(x)\)，即 \(f\) 是偶函数.

令 \(u(x)=\sqrt{1+x^{2}}-x\). 由于 \(\sqrt{1+x^{2}}>|x|\)，所以 \(u(x)>0\)，且 \(u(-x)=\sqrt{1+x^{2}}+x\). 于是 \(u(x)u(-x)=\left(\sqrt{1+x^{2}}-x\right)\left(\sqrt{1+x^{2}}+x\right)=1\). 因此 \(g(-x)=\lg u(-x)=\lg\frac1{u(x)}=-\lg u(x)=-g(x)\)，即 \(g\) 是奇函数. 故选 B.
\end{solution}

\begin{problem}
若等差数列的第一，二，三项依次是 \(\frac1{x+1}\)，\(\frac5{6x}\)，\(\frac12\)，那么，这个等差数列的第 \(101\) 项的值是（\quad）

\choices
{\(50\frac13\)}
{\(13\frac23\)}
{\(24\)}
{\(8\frac23\)}
\end{problem}

\begin{solution}
等差数列满足 \(2a_{2}=a_{1}+a_{3}\)，所以 \(\frac{5}{3x}=\frac1{x+1}+\frac12\). 由于 \(x\neq0\)，\(-1\)，通分得 \(3x^{2}-x-10=0\)，即 \((3x+5)(x-2)=0\).

当 \(x=2\) 时，\(a_{1}=\frac13\)，公差 \(d=\frac{5}{12}-\frac13=\frac1{12}\)，故 \(a_{101}=\frac13+100\cdot\frac1{12}=\frac{26}{3}=8\frac23\)，对应选项 D.

当 \(x=-\frac53\) 时，\(a_{1}=-\frac32\)，\(a_{2}=-\frac12\)，\(a_{3}=\frac12\)，三项同样构成等差数列，公差为 \(1\)，从而 \(a_{101}=-\frac32+100=\frac{197}{2}\)，但该值不在选项中. 因此，严格按原题面不能唯一确定选项；若题面本意是补充 \(x>0\) 条件，则取 \(x=2\)，答案为 D.
\end{solution}

\section{填空题}

\begin{problem}
若 \(z=\sqrt2+\i\)，则 \(z^{2}\) 的共轭复数的代数形式是\fillinblank{}.
\end{problem}

\begin{answer}\(1-2\sqrt2\i\)\end{answer}

\begin{solution}
先平方得 \(z^{2}=(\sqrt2+\i)^{2}=2+2\sqrt2\i+\i^{2}=1+2\sqrt2\i\). 所以 \(z^{2}\) 的共轭复数为 \(1-2\sqrt2\i\).
\end{solution}

\begin{problem}
已知椭圆的两个焦点是 \(F_{1}(-2,0)\) 和 \(F_{2}(2,0)\)，并且点 \(A(0,2)\) 在该椭圆上，那么这个椭圆的标准方程是\fillinblank{}.
\end{problem}

\begin{answer}\(\displaystyle\frac{x^{2}}8+\frac{y^{2}}4=1\)\end{answer}

\begin{solution}
由 \(A(0,2)\) 到两个焦点的距离相等，且
\(AF_{1}=AF_{2}=\sqrt{(0+2)^{2}+2^{2}}=2\sqrt2\). 椭圆的焦点定义给出 \(2a=AF_{1}+AF_{2}=4\sqrt2\)，所以 \(a=2\sqrt2\)，\(a^{2}=8\). 又焦半距 \(c=2\)，故 \(b^{2}=a^{2}-c^{2}=8-4=4\)，标准方程为 \(\frac{x^{2}}8+\frac{y^{2}}4=1\).
\end{solution}

\begin{problem}
\(\displaystyle\lim_{n\to\infty}\frac{4n^{3}-1}{n(2n^{2}+1)}\) 的值是\fillinblank{}.
\end{problem}

\begin{answer}\(2\)\end{answer}

\begin{solution}
分子、分母同除以 \(n^{3}\)，得
\(\displaystyle\lim_{n\to\infty}\frac{4-\frac1{n^{3}}}{2+\frac1{n^{2}}}=\frac42=2\).
\end{solution}

\begin{problem}
设 \((1-x)^{7}=a_{0}+a_{1}x+a_{2}x^{2}+a_{3}x^{3}+a_{4}x^{4}+a_{5}x^{5}+a_{6}x^{6}+a_{7}x^{7}\)，则 \(a_{0}+a_{2}+a_{4}+a_{6}\) 的值是\fillinblank{}.
\end{problem}

\begin{answer}\(64\)\end{answer}

\begin{solution}
令 \(E=a_{0}+a_{2}+a_{4}+a_{6}\)，\(O=a_{1}+a_{3}+a_{5}+a_{7}\). 取 \(x=1\)，得 \(E+O=(1-1)^{7}=0\)；取 \(x=-1\)，得 \(E-O=(1+1)^{7}=128\). 两式相加得 \(2E=128\)，所以 \(E=64\).
\end{solution}

\begin{problem}
函数 \(y=\frac{2x}{1+x^{2}}\)（\(x\in\R\)）的值域是\fillinblank{}.
\end{problem}

\begin{answer}\([-1,1]\)\end{answer}

\begin{solution}
因为 \(1+x^{2}>0\)，且 \((x-1)^{2}\geqslant0\)，所以 \(1+x^{2}\geqslant2x\)，从而 \(\frac{2x}{1+x^{2}}\leqslant1\)；又因为 \((x+1)^{2}\geqslant0\)，所以 \(1+x^{2}\geqslant-2x\)，从而 \(\frac{2x}{1+x^{2}}\geqslant-1\). 当 \(x=1\) 和 \(x=-1\) 时分别取得 \(1\) 和 \(-1\). 函数 \(y=\frac{2x}{1+x^{2}}\) 在 \(\R\) 上连续，由介值定理可知 \([-1,1]\) 内的每个值都能取得，所以值域为 \([-1,1]\).
\end{solution}

\section{解答题}

\begin{problem}
设正数 \(a\)，\(b\)，\(c\) 满足 \(a^{2}+b^{2}=c^{2}\).
\begin{enumerate}
\item 求证：\(\log_{2}\left(1+\frac{b+c}{a}\right)+\log_{2}\left(1+\frac{a-c}{b}\right)=1\).
\item 又设 \(\log_{4}\left(1+\frac{b+c}{a}\right)=1\)，\(\log_{8}(a+b-c)=\frac23\). 求 \(a\)，\(b\)，和 \(c\) 的值.
\end{enumerate}
\end{problem}

\begin{solution}
\begin{enumerate}
\item
由 \(a^{2}+b^{2}=c^{2}\) 且 \(a\)，\(b\)，\(c>0\)，有 \(c<a+b\)，所以两个对数的真数均为正. 于是
\(\log_{2}\left(1+\frac{b+c}{a}\right)+\log_{2}\left(1+\frac{a-c}{b}\right)=\log_{2}\left(\frac{a+b+c}{a}\cdot\frac{a+b-c}{b}\right)\)
\(=\log_{2}\frac{(a+b)^{2}-c^{2}}{ab}=\log_{2}\frac{a^{2}+2ab+b^{2}-a^{2}-b^{2}}{ab}=\log_{2}2=1\).

\item
由 \(\log_{4}\left(1+\frac{b+c}{a}\right)=1\)，得 \(\frac{a+b+c}{a}=4\)，即 \(a+b+c=4a\). 又由 \(\log_{8}(a+b-c)=\frac23\)，得 \(a+b-c=8^{2/3}=4\).

第一问中的恒等式说明两个对数之和为 \(1\)，而第一个对数为 \(2\)，所以 \(\log_{2}\frac{a+b-c}{b}=-1\)，即 \(a+b-c=\frac b2\). 结合 \(a+b-c=4\)，得 \(b=8\). 再由 \(a+b-c=4\) 得 \(c=a+4\)，代入 \(a+b+c=4a\)，得 \(a+8+(a+4)=4a\)，所以 \(a=6\)，\(c=10\).
\end{enumerate}
\end{solution}

\begin{problem}
在三棱台 \(A_{1}B_{1}C_{1}-ABC\) 中，侧棱 \(B_{1}B\perp\) 底面 \(ABC\)，且 \(\angle ABC=\angle AA_{1}C=\frac\pi2\)，\(AB=2A_{1}B_{1}=2\mathrm{~cm}\).
\begin{enumerate}
\item 求证：\(BC\perp A_{1}B\)，\(BC\perp A_{1}A\)，\(A_{1}A\perp A_{1}B\).
\item 求异面直线 \(A_{1}A\) 和 \(BC\) 的距离.
\end{enumerate}
\bitmapfigure[max width=0.46\linewidth]{img/1987/guangdong_liberal/q7_fig1.png}
\end{problem}

\begin{solution}
\begin{enumerate}
\item
因为 \(B_{1}B\perp\) 平面 \(ABC\)，所以 \(BC\perp B_{1}B\)；又由 \(\angle ABC=\frac\pi2\)，得 \(BC\perp AB\). 直线 \(AB\) 和 \(BB_{1}\) 是平面 \(ABB_{1}A_{1}\) 内相交的两条直线，因此 \(BC\perp\) 平面 \(ABB_{1}A_{1}\)，从而 \(BC\perp A_{1}B\)，\(BC\perp A_{1}A\).
又由 \(\angle AA_{1}C=\frac\pi2\)，得 \(AA_{1}\perp A_{1}C\). 结合 \(AA_{1}\perp BC\)，且 \(A_{1}C\) 与 \(BC\) 在点 \(C\) 相交，知 \(AA_{1}\perp\) 平面 \(A_{1}BC\)，所以 \(AA_{1}\perp A_{1}B\).
\item
由 \(AB=2A_{1}B_{1}=2\mathrm{~cm}\)，得 \(A_{1}B_{1}=1\mathrm{~cm}\). 由棱台性质 \(AB\parallel A_{1}B_{1}\)，且侧棱垂直底面、两底面平行，得 \(B_{1}B\perp A_{1}B_{1}\). 所以 \(\angle B_{1}A_{1}B=\angle ABA_{1}\)，又 \(A_{1}A\perp A_{1}B\)，故直角三角形 \(\triangle A_{1}B_{1}B\) 与 \(\triangle BA_{1}A\) 相似，得到 \(\frac{A_{1}B_{1}}{A_{1}B}=\frac{A_{1}B}{AB}\). 于是 \((A_{1}B)^{2}=A_{1}B_{1}\cdot AB=1\cdot2=2\)，所以 \(A_{1}B=\sqrt2\mathrm{~cm}\).
由第（1）问，\(A_{1}B\perp A_{1}A\) 且 \(A_{1}B\perp BC\)，故 \(A_{1}B\) 是两条异面直线的公垂线，所求距离为 \(\sqrt2\mathrm{~cm}\).
\end{enumerate}
\end{solution}

\begin{problem}
直线 \(l\) 通过抛物线 \(y^{2}=2px\)（\(p\neq0\)）的焦点，并且与这抛物线相交于 \(A(x_{1},y_{1})\) 和 \(B(x_{2},y_{2})\) 两点.
\begin{enumerate}
\item 求证：\(y_{1}y_{2}=-p^{2}\).
\item 点 \(C\) 在这抛物线的准线上，且 \(AC\) 平行于 \(x\) 轴，求证：\(B\)，\(C\) 和这抛物线的顶点三点共线.
\end{enumerate}
\end{problem}

\begin{solution}
\FigureLayoutDeclare{img/1987/guangdong_liberal/q33_solution_fig1.png}{8.0cm}{0bp 0bp 0bp 8.801bp}
\solutionfigure{\bitmapfigure[max width=0.42\linewidth]{img/1987/guangdong_liberal/q33_solution_fig1.png}}
\begin{enumerate}
\item
抛物线的焦点为 \(F\left(\frac p2,0\right)\). 若 \(l\perp x\) 轴，则 \(x_{1}=x_{2}=\frac p2\)，代入抛物线得 \(y_{1}^{2}=y_{2}^{2}=p^{2}\). 由于 \(A\)，\(B\) 是两个交点，\(y_{1}\) 与 \(y_{2}\) 取相反数，故 \(y_{1}y_{2}=-p^{2}\).

若 \(l\) 不垂直于 \(x\) 轴，由于 \(A\)，\(B\) 是两个不同的交点，\(l\) 不可能是 \(x\) 轴（\(x\) 轴与该抛物线只有顶点一个交点），故可设 \(l:y=k\left(x-\frac p2\right)\)，其中 \(k\neq0\). 于是 \(x=\frac yk+\frac p2\)，代入抛物线得 \(y^{2}-\frac{2p}{k}y-p^{2}=0\). 其两个根正是 \(y_{1}\)，\(y_{2}\)，由韦达定理得 \(y_{1}y_{2}=-p^{2}\).
\item
抛物线的准线为 \(x=-\frac p2\). 由 \(AC\parallel x\) 轴，得 \(C\left(-\frac p2,y_{1}\right)\). 由 \(y_{1}y_{2}=-p^{2}\) 知 \(y_{2}\neq0\)，又 \(x_{2}=\frac{y_{2}^{2}}{2p}\)，所以
\(\frac{y_{2}}{x_{2}}=\frac{2p}{y_{2}}=-\frac{2y_{1}}{p}=\frac{y_{1}}{-p/2}\).
这说明直线 \(OB\) 与 \(OC\) 的斜率相等，其中 \(O(0,0)\) 是抛物线的顶点. 因此 \(B\)，\(C\) 和 \(O\) 三点共线.
\end{enumerate}
\end{solution}

\begin{problem}
设 \(|\theta|<\frac\pi2\)，数列 \(\{a_n\}\) 的通项是 \(a_n=\sin\frac{n\pi}{2}\cdot\tan^{n}\theta\)，记 \(S_{2n}=a_{1}+a_{2}+\dots+a_{2n-1}+a_{2n}\).
\begin{enumerate}
\item 求证：对任何自然数 \(n\)，\(S_{2n}=\frac12(\sin2\theta)[1+(-1)^{n+1}\tan^{2n}\theta]\).
\item 求证：当 \(n\geqslant2\) 时，\(a_n^{2}+a_{n-1}a_{n+1}+(-1)^n(a_{1}^{n})^{2}=0\).
\end{enumerate}
\end{problem}

\begin{solution}
\begin{enumerate}
\item
令 \(t=\tan\theta\). 因为 \(\sin\frac{2k\pi}{2}=0\)，\(\sin\frac{(2k-1)\pi}{2}=(-1)^{k+1}\)，所以 \(a_{2k}=0\)，\(a_{2k-1}=(-1)^{k+1}t^{2k-1}\). 因此
\(S_{2n}=t-t^{3}+\cdots+(-1)^{n+1}t^{2n-1}=\frac{t\left[1-(-t^{2})^{n}\right]}{1+t^{2}}=\frac{t\left[1+(-1)^{n+1}t^{2n}\right]}{1+t^{2}}=\frac12(\sin2\theta)\left[1+(-1)^{n+1}\tan^{2n}\theta\right]\).
\item
注意 \(a_{1}=t\). 当 \(n=2k\) 时，\(a_{n}=0\)，且
\(a_{n-1}a_{n+1}=a_{2k-1}a_{2k+1}=(-1)^{k+1}t^{2k-1}\cdot(-1)^{k+2}t^{2k+1}=-t^{4k}\)，
\((-1)^{n}(a_{1}^{n})^{2}=t^{4k}\). 三项之和为 \(0-t^{4k}+t^{4k}=0\).
当 \(n=2k-1\) 时，\(a_{n-1}a_{n+1}=0\)，且 \(a_{n}^{2}=t^{4k-2}\)，\((-1)^{n}(a_{1}^{n})^{2}=-t^{4k-2}\). 三项之和为 \(t^{4k-2}+0-t^{4k-2}=0\). 所以结论成立.
\end{enumerate}
\end{solution}
